\documentclass{article}
\usepackage[T1]{fontenc}
\usepackage{lmodern}
\usepackage{graphicx}
\usepackage{amsmath,amssymb}
\usepackage[a4paper,margin=2cm]{geometry}
\usepackage[absolute,overlay]{textpos}
\setlength{\TPHorizModule}{1mm}
\setlength{\TPVertModule}{1mm}
\usepackage[indonesian]{babel}
\usepackage{hyperref}
\setlength{\emergencystretch}{2em}
% unit: Halaman 19

\begin{document}

% === Halaman 19 (llm) ===

\section*{Enkripsi}

\begin{itemize}
    \item Tinjau $RC5$ dengan ukuran blok 64 bit dan jumlah putaran $r$.
    \item Enkripsi menggunakan kunci putaran $S_0, S_1, \dots, S_{2r+2}$ yang masing-masing panjangnya 32-bit.
    \item Dua kunci putaran digunakan untuk setiap putaran $i = 1, 2, \dots, r$ dan dua buah kunci putaran tambahan sebelum putaran pertama jadi seluruhnya ada $2r + 2$ buah kunci internal.
    \item Untuk melakukan enkripsi, mula-mula blok plainteks dibagi menjadi 2 bagian, $A$ dan $B$, yang masing-masing panjangnya 32 bit. Kemudian masing-masing bagian dijumlahkan (dalam modulo $2^{32}$) dengan $S_0$ dan $S_1$:
    \[
    \begin{aligned}
    A \leftarrow A + S[0] \\
    B \leftarrow B + S[1]
    \end{aligned}
    \]
\end{itemize}

\end{document}
